% TikZ 3D schematic of the Newell cell-pair geometry,
% used in docs/theory.tex (Section 5.3.3, Level 2).
%
% Standalone document: compile with `pdflatex newell_geometry.tex`
% to produce newell_geometry.pdf in the same directory.
\documentclass[tikz, border=6pt]{standalone}
\usepackage{tikz-3dplot}
\usepackage{amsmath}
\usetikzlibrary{arrows.meta, calc}
% Theta = polar angle from +z; phi = azimuth from +x in xy-plane.
\tdplotsetmaincoords{70}{70}

\begin{document}
\begin{tikzpicture}[
    tdplot_main_coords,
    scale=1.8,
    every node/.style={font=\small},
    surface/.style={opacity=0.55, draw=black, line width=0.4pt},
    wire/.style={black, line width=0.6pt},
    Rarrow/.style={-{Stealth[length=4pt]}, violet, line width=0.9pt},
    Marrow/.style={-{Stealth[length=4pt]}, black, line width=1.1pt},
    axarr/.style={-{Stealth[length=3pt]}, gray, line width=0.5pt},
]

% --- Geometry constants ------------------------------------------
\def\hsx{0.50} % half a along x
\def\hsy{0.50} % half a along y
\def\hsz{0.30} % half t_Co along z
\def\Rx{2.60}  % source-to-dest displacement, x
\def\Ry{1.40}  % displacement, y
\def\Rz{0.50}  % displacement, z

% --- Source cell: wireframe edges, all 12 of them ----------------
\draw[wire] (-\hsx,-\hsy,-\hsz) -- (\hsx,-\hsy,-\hsz)
            -- (\hsx,\hsy,-\hsz) -- (-\hsx,\hsy,-\hsz) -- cycle;
\draw[wire] (-\hsx,-\hsy,\hsz) -- (\hsx,-\hsy,\hsz)
            -- (\hsx,\hsy,\hsz) -- (-\hsx,\hsy,\hsz) -- cycle;
\draw[wire] (-\hsx,-\hsy,-\hsz) -- (-\hsx,-\hsy,\hsz);
\draw[wire] (\hsx,-\hsy,-\hsz) -- (\hsx,-\hsy,\hsz);
\draw[wire] (\hsx,\hsy,-\hsz) -- (\hsx,\hsy,\hsz);
\draw[wire] (-\hsx,\hsy,-\hsz) -- (-\hsx,\hsy,\hsz);

% --- Source +u face (x = +hsx) shaded red (+sigma) --------------
\fill[red!55, surface]
    (\hsx,-\hsy,-\hsz) -- (\hsx,\hsy,-\hsz)
    -- (\hsx,\hsy,\hsz) -- (\hsx,-\hsy,\hsz) -- cycle;
% --- Source -u face (x = -hsx) shaded blue (-sigma) --------------
\fill[blue!55, surface]
    (-\hsx,-\hsy,-\hsz) -- (-\hsx,\hsy,-\hsz)
    -- (-\hsx,\hsy,\hsz) -- (-\hsx,-\hsy,\hsz) -- cycle;

% --- Source-surface integration nodes (3 x 3 on +u face) ---------
\foreach \yy in {-0.32, 0, 0.32}{
  \foreach \zz in {-0.18, 0, 0.18}{
    \filldraw[red!70!black]
        (\hsx, \yy, \zz) circle (1.0pt);
  }
}

% --- Magnetization arrow inside source: M = M_u u-hat ------------
\draw[Marrow] (-0.30, 0, 0.05) -- (0.55, 0, 0.05);

% --- Destination cell: wireframe (offset by R) --------------------
\draw[wire] (\Rx-\hsx,\Ry-\hsy,\Rz-\hsz)
       -- (\Rx+\hsx,\Ry-\hsy,\Rz-\hsz)
       -- (\Rx+\hsx,\Ry+\hsy,\Rz-\hsz)
       -- (\Rx-\hsx,\Ry+\hsy,\Rz-\hsz) -- cycle;
\draw[wire] (\Rx-\hsx,\Ry-\hsy,\Rz+\hsz)
       -- (\Rx+\hsx,\Ry-\hsy,\Rz+\hsz)
       -- (\Rx+\hsx,\Ry+\hsy,\Rz+\hsz)
       -- (\Rx-\hsx,\Ry+\hsy,\Rz+\hsz) -- cycle;
\draw[wire] (\Rx-\hsx,\Ry-\hsy,\Rz-\hsz)
       -- (\Rx-\hsx,\Ry-\hsy,\Rz+\hsz);
\draw[wire] (\Rx+\hsx,\Ry-\hsy,\Rz-\hsz)
       -- (\Rx+\hsx,\Ry-\hsy,\Rz+\hsz);
\draw[wire] (\Rx+\hsx,\Ry+\hsy,\Rz-\hsz)
       -- (\Rx+\hsx,\Ry+\hsy,\Rz+\hsz);
\draw[wire] (\Rx-\hsx,\Ry+\hsy,\Rz-\hsz)
       -- (\Rx-\hsx,\Ry+\hsy,\Rz+\hsz);

% Light fill for the dest cell so it reads as a body
\fill[yellow!30, opacity=0.20]
    (\Rx-\hsx,\Ry-\hsy,\Rz-\hsz)
    -- (\Rx+\hsx,\Ry-\hsy,\Rz-\hsz)
    -- (\Rx+\hsx,\Ry+\hsy,\Rz-\hsz)
    -- (\Rx-\hsx,\Ry+\hsy,\Rz-\hsz) -- cycle;
\fill[yellow!30, opacity=0.20]
    (\Rx-\hsx,\Ry-\hsy,\Rz+\hsz)
    -- (\Rx+\hsx,\Ry-\hsy,\Rz+\hsz)
    -- (\Rx+\hsx,\Ry+\hsy,\Rz+\hsz)
    -- (\Rx-\hsx,\Ry+\hsy,\Rz+\hsz) -- cycle;

% --- Destination volume integration nodes (3 x 3 x 3) -------------
\foreach \xx in {-0.32, 0, 0.32}{
  \foreach \yy in {-0.32, 0, 0.32}{
    \foreach \zz in {-0.18, 0, 0.18}{
      \filldraw[olive!80!black]
         (\Rx+\xx, \Ry+\yy, \Rz+\zz) circle (0.7pt);
    }
  }
}

% --- Displacement vector R: source center -> destination center ---
\draw[Rarrow] (0,0,0) -- (\Rx,\Ry,\Rz);

% --- Cartesian-axis triad at the lower-left ----------------------
\def\ax{-1.2}
\def\ay{-1.2}
\def\az{-0.65}
\draw[axarr] (\ax,\ay,\az) -- (\ax+0.55,\ay,\az);
\draw[axarr] (\ax,\ay,\az) -- (\ax,\ay+0.55,\az);
\draw[axarr] (\ax,\ay,\az) -- (\ax,\ay,\az+0.45);

% --- Text labels --------------------------------------------------
\node[above=10pt, font=\small\bfseries] at (0, 0, \hsz)
  {source cell};
\node[above=10pt, font=\small\bfseries]
  at (\Rx, \Ry, \Rz+\hsz) {destination cell};

\node[right=2pt, black] at (0.30, 0, -0.22)
  {$\mathbf{M} = M_u\,\hat{u}$};
\node[below=2pt, red!60!black] at (\hsx, 0, -\hsz)
  {\small $+\sigma = +M_u$};
\node[below=2pt, blue!60!black] at (-\hsx, 0, -\hsz)
  {\small $-\sigma = -M_u$};

\node[above=2pt, violet, font=\bfseries]
  at (\Rx*0.40, \Ry*0.40, \Rz*0.50 + 0.18) {$\mathbf{R}$};

\node[gray, right] at (\ax+0.55, \ay, \az)
  {$x \equiv u$};
\node[gray, above right] at (\ax, \ay+0.55, \az)
  {$y \equiv v$};
\node[gray, left] at (\ax, \ay, \az+0.45)
  {$z \equiv w$};

% --- Dimension labels on destination cell -------------------------
\node[font=\footnotesize, anchor=west]
  at (\Rx+\hsx+0.06, \Ry+\hsy, \Rz-\hsz)
  {$a$};
\node[font=\footnotesize, anchor=west]
  at (\Rx-\hsx, \Ry-\hsy-0.10, \Rz-\hsz-0.18)
  {$a$};
\node[font=\footnotesize, anchor=west]
  at (\Rx-\hsx-0.10, \Ry-\hsy, \Rz)
  {$t_{\text{Co}}$};

% (legend omitted; rely on the LaTeX caption instead)

\end{tikzpicture}
\end{document}
